%%%%%
%%%%%  Naudokite LUALATEX, ne LATEX.
%%%%%
%%%%
\documentclass[]{VUMIFTemplateClass}

\usepackage{indentfirst}
\usepackage{amsmath, amsthm, amssymb, amsfonts}
\usepackage{mathtools}
\usepackage{physics}
\usepackage{graphicx}
\usepackage{verbatim}
\usepackage[hidelinks]{hyperref}
\usepackage{color,algorithm,algorithmic}
\usepackage[nottoc]{tocbibind}
\usepackage{tocloft}

\usepackage{titlesec}
\newcommand{\sectionbreak}{\clearpage}

\makeatletter
\renewcommand{\fnum@algorithm}{\thealgorithm}
\makeatother
\renewcommand\thealgorithm{\arabic{algorithm} algoritmas}

\usepackage{biblatex}
\bibliography{bibliografija}
%% norint pakeisti bibliografijos šaltinių numeravimą (skaitiniu arba raidiniu), pakeitimus atlikti VUMIFTemplateClass.cls 150 eilutėje

% Author's MACROS
\newcommand{\EE}{\mathbb{E}\,} % Mean
\newcommand{\ee}{{\mathrm e}}  % nice exponent
\newcommand{\RR}{\mathbb{R}}


\studijuprograma{Programų sistemų} %Studijų programą įrašyti kilmininko linksniu (pavyzdžiui – Programų sistemų, Finansų ir draudimų matematikos ir t. t.)
\darbotipas{Bakalauro baigiamojo darbo planas} % Bakalauro baigiamasis darbas arba magistro baigiamasis darbas
\darbopavadinimas{Maišymo proceso modeliavimas YAG reakcijose}
\darbopavadinimasantras{Modelling the mixing process in YAG reactions}
\autorius{Arnas Vaicekauskas}

%Autorių gali būti ir daugiau, tuo atveju, kiekvienas autorius rašomas iš naujos eilutės, ir pridedamas titulinis.tex arba dvigubasTitulinis.tex dokumentuose
%\antrasautorius{Vardas Pavardė} %Jei toks yra, kitu atveju ištrinti

\vadovas{Asist. Dr. Rokas Astrauskas}
\recenzentas{pedagoginis/mokslinis vardas Vardas Pavardė} %Jei toks yra žinomas, kitu atveju ištrinti

\begin{document}
\selectlanguage{lithuanian}

\onehalfspacing
\begin{titlepage}
\vskip 20pt
\begin{center}
\includegraphics[scale=0.55]{images/MIF.png}
\end{center}

\makeatletter

\vskip 20pt
\centerline{\bf \large \textbf{VILNIAUS UNIVERSITETAS}}
\vskip 10pt
\centerline{\large \textbf{MATEMATIKOS IR INFORMATIKOS FAKULTETAS}}
\vskip 10pt
\centerline{\large \textbf{\MakeUppercase{\@studijuprograma \space studijų programa}}}

\vskip 80pt
\centerline{\Large \@darbotipas}
\vskip 20pt
\begin{center}
    {\bf \LARGE \@darbopavadinimas}
\end{center}
\begin{center}
    {\bf \Large \@darbopavadinimasantras}
\end{center}
\vskip 80pt

\centering{\Large \@autorius}
\@ifundefined{@antrasautorius}{}
{
\vskip 10pt
\centering{\Large \@antrasautorius}
}
\vskip 20pt

\centering{
    \begin{tabular}{rcp{.7\textwidth}}
        {\Large Darbo vadovas} & {\Large :} & {\Large \@vadovas}\\[10pt]
        \@ifundefined{@moksliniskonsultantas}{}
            {
                {\Large Mokslinis konsultantas} & {\Large :} & {\Large \@moksliniskonsultantas}\\[10pt]
            }
        \@ifundefined{@recenzentas}{}
            {
                {\Large Recenzentas} & {\Large :} & {\Large \@recenzentas}\\[10pt]
            }
    \end{tabular}}


\vskip 110pt

\centerline{\large \textbf{Vilnius}}
\centerline{\large \textbf{\the\year{}}}

\makeatother

\newpage
\end{titlepage}
%\newgeometry{top=2cm,bottom=2cm,right=2cm,left=3cm}
\setcounter{page}{2}


%% Padėkų skyrius
\section{Darbo planas}

\subsection{Tyrimo objektas ir aktualumas}

% Itrio aliuminio granato (YAG) kristalai, legiruoti neodimio jonais, yra viena populiariausių medžiagų lazerių aktyviosioms terpėms, naudojamoms pramoninėje gamyboje ir medicinoje. Chemiškai paprasčiausias ir pramoninei gamybai labiausiai pritaikytas šios medžiagos sintezės metodas yra kietafazė reakcija, kurios spartinimui taikomas reagentų išmaišymas. Šiame darbe analizuojamas YAG kietafazės reakcijos modelis, į kurį įtrauktas išmaišymo mechanizmas.

% outline:
% YAG medžiaga
% optinės savybės
% pritaikymas rinkoje (trumpai, gal kažką pacituoti)
% sintezės būdai, kietafazės reakcijos paprastumas bei pritaikomumas pramonėje (turėjau lygtais čia source vieną)
% modeliavimo privalumai (reakcijos mechanizmo aiškinimasis, kaina, laikas, žmogiškieji ištekliai)
% modeliavimo būdai, FEM, BEM ir naudojamas FDM (finite difference method)
% konkretūs finite difference metodai:
%   explicit FTCS, jo paprastumas ir trūkumai
%   implicit ADI, stiprybės ir trūkumai, dabartiniai taikymai
% didesnės erdvės? privesti prie tikslo?

% Itrio aliuminio granatas arba kitaip - YAG, yra kristalinis junginys, kuris pasižymi išskirtinėmis optinėmis savybėmis. Dėl šių savybiu, YAG kristalai, legiruoti su neodimio jonais, yra plačiai naudojami kaip lazerių aktyvioji terpė. Šios medžiagos lazeriai yra plačiai naudojami odontologijoje \cite{valentiUseErYAG2021}, pramoninė gamyboje \cite{dubeyExperimentalStudyNd2008} bei aibėje kitų sričių.

% Ši medžiaga gali būti sintezuota keliais skirtingais būdais - kietafazė reakcija, solvoterminiu procesu, nusodinimu, purškiamo aerozolio terminiu skilimu bei zolio-gelio procesu. Daugelis iš išvardintų metodų yra chemiškai sudėtingi ir privalo būti atliekami laboratorijoje. Paprasčiausias ir geriausiai pramoninei gamybai pritaikomas metodas sintezuoti šią medžiaga yra kietafazė reakcija \cite{zhangNovelSynthesisYAG2005}. Šios reakcijos metu iš Itrio ir aliuminio oksidai reaguoja aukštoje temperatūroje, kurioje ir susidaro YAG kristalai. Tiesa, praktikoje šią reakcija galima spartinti reagentus periodiškai išmaišant.

% Kompiuterinis fizinių bei cheminių procesų modeliavimas yra ypač paplitusi tyrimo metodika, kuri leidžia giliau suprasti ir tiksliau nuspėti tokio pobūdžio procesų eigą ir jų rezultatus. Tai galioja ir šiuo atveju - laboratorijoje atliekant YAG sintezės reakciją, chemikai turi ribotas galimybes ištirti reakcijos metu veikiančius mechanizmus, todėl šiame darbe modeliuosime YAG sintezės reakcijos modelį kaip reakcijos-difuzijos sistemą. Reakcijos-difuzijos sistemos yra senas, tačiau vis dar plačiai taikomas metodas modeliuoti tokio tipo procesus. Matematinį YAG reakcijos modelį pasiūlė Ivanauskas et. al. \cite{ivanauskasModellingSolidState2005}. Šiame ir susijusiuose straipsniuose \cite{ivanauskasComputationalModellingYAG2009,mackeviciusCloserLookComputer2012} eksperimentiniu būdu yra išmatuotos fizinės modelio konstantos todėl kompiuterinio modelio rezultatai gali būti reikšmingai matuojami ir lyginami.

% Šią reakcija nusako trejetas netiesinių parabolinių dalinių diferencialinių lygčių, kurios sudaro reakcijos-difuzijos sistemą. Tokio tipo sistemoms uždara sprendinio forma neegzistuoja todėl jų sprendimui yra taikomi skaitiniai metodai. Spręsti dalinių diferencialinių lygčių sistemas dažniausiai yra naudojami tokie metodai kaip baigtinių elementų metodas (\textit{angl. finite element method, FEM}) - <neilgas sakinukas apie šio metodo esmę>. Taip pat gali būti taikomas ribinių elementų metodas (\textit{angl. boundary element method, BEM}), kuris <trumpas sakinukas apie tai kaip veikia ir kada naudojamas>. Šiuo atveju, spręsti šiai sistemai naudosime du baigtinių skirtumų metodus (\textit{angl. finite difference method, FDM}). Pirmasis metodas yra išreikštinis FTCS metodas arba kitaip žinomas kaip Oilerio integracija - šis metodas nėra techniškai sudėtingas, tačiau yra sąlyginai stabilus ir nėra ypatingai tikslus. Antrasis metodas yra žinomas kaip neišreikštinis kintamosios krypties metodas (\textit{angl. alternating direction implicit, ADI}) - šis metodas yra techniškai sudėtingas, tačiau nesąlyginiai stabilus bei turi didesnį tikslumą \cite{NumericalSolutionParabolic}. Nors ADI metodas yra senas, tačiau vis dar plačiai taikomas \cite{gaidamauskaiteComparisonFiniteDifference2007}.

Itrio aliuminio granatas (YAG) yra kristalinis junginys, pasižymintis išskirtinėmis optinėmis savybėmis. Dėl šių savybių YAG kristalai, legiruoti neodimio jonais, plačiai naudojami kaip lazerių aktyvioji terpė. Tokie lazeriai taikomi įvairiose srityse, įskaitant odontologiją \cite{valentiUseErYAG2021}, pramoninę gamybą \cite{dubeyExperimentalStudyNd2008} ir daugelį kitų.

YAG gali būti sintezuojamas keliais skirtingais būdais, įskaitant kietafazę reakciją, solvoterminį procesą, nusodinimą, purškiamo aerozolio terminį skilimą ir zolio-gelio metodą. Daugelis šių metodų yra chemiškai sudėtingi ir reikalauja specializuotos laboratorinės įrangos. Tarp jų kietafazė reakcija išsiskiria kaip paprasčiausias ir pramoninei gamybai tinkamiausias sintezės būdas \cite{zhangNovelSynthesisYAG2005}. Šios reakcijos metu itrio ir aliuminio oksidai reaguoja aukštoje temperatūroje, susidarant YAG kristalams. Praktikoje šios reakcijos greitį galima padidinti periodiškai maišant reagentus.

Kompiuterinis fizinių bei cheminių procesų modeliavimas yra plačiai taikomas tyrimo metodas, leidžiantis giliau suprasti ir tiksliau prognozuoti tokių procesų eigą bei jų rezultatus. Šis metodas itin naudingas YAG sintezės reakcijų analizėje, kadangi laboratoriniai eksperimentai suteikia ribotas galimybes detaliai ištirti vykstančius mechanizmus. Todėl šiame darbe modeliuosime YAG sintezės reakciją kaip reakcijos-difuzijos sistemą. Nors šis metodas taikomas jau ilgą laiką, jis išlieka aktualus ir šiandien. Matematinį YAG reakcijos modelį pasiūlė Ivanauskas et al. \cite{ivanauskasModellingSolidState2005}, o susijusiuose darbuose \cite{ivanauskasComputationalModellingYAG2009,mackeviciusCloserLookComputer2012} eksperimentiniu būdu nustatytos fizinės modelio konstantos, todėl skaičiavimų rezultatus galima tiksliai palyginti su eksperimentiniais duomenimis. Minėtuose straipsniuose dėmesys skiriamas pačiai YAG sintezės reakcijai ir koeficientų nustatymui ir dėl to reakcija yra modeliuojama labai mažoje fizinėje erdvėje, kurioje telpa vos kelios mikrodalelės. Šiame darbe yra tiriamas maišymo mechanizmo poveikis reakcijai, kuris gali kisti priklausomai nuo fizinės erdvės dydžio, todėl yra reikšminga modeliuoti įvairaus dydžio erdves.

Šią reakciją aprašo trijų netiesinių parabolinių dalinių diferencialinių lygčių sistema, sudaranti reakcijos-difuzijos modelį. Tokiose sistemose analitinis sprendinys paprastai neegzistuoja, todėl jų sprendimui taikomi skaitiniai metodai. Vienas iš dažniausiai naudojamų metodų yra baigtinių elementų metodas (\textit{angl. finite element method, FEM}) - jis leidžia suskaidyti nagrinėjamą sritį į mažesnius elementus ir spręsti diferencialines lygtis kiekviename iš jų. Taip pat gali būti taikomas ribinių elementų metodas (\textit{angl. boundary element method, BEM}), kuris naudoja tik ribinių sąlygų informaciją, todėl yra efektyvus sprendžiant uždavinius su sudėtingomis geometrijomis.

Šiame darbe reakcijos-difuzijos sistemai spręsti taikysime du baigtinių skirtumų metodus (\textit{angl. finite difference method, FDM}). Pirmasis - išreikštinis FTCS metodas (\textit{angl. Forward-Time Centered-Space}), dar žinomas kaip Oilerio integracija. Šis metodas pasižymi paprasta implementacija, tačiau yra sąlyginai stabilus ir nėra itin tikslus. Antrasis metodas - neišreikštinis kintamosios krypties metodas (\textit{angl. alternating direction implicit, ADI}), kuris yra techniškai sudėtingesnis, tačiau nesąlyginai stabilus ir užtikrina didesnį tikslumą \cite{doi:10.1137/0103003}. Nepaisant to, kad ADI metodas sukurtas jau seniai, jis vis dar plačiai naudojamas \cite{gaidamauskaiteComparisonFiniteDifference2007}. Šie skaitiniai metodai detaliai aprašyti literatūroje \cite{pressNumericalRecipes3rd2007,levequeFiniteDifferenceMethods2007}. 

% Šio tyrimo tikslas - ištirti maišymo poveikį reakcijai įvairaus dydžio erdvėse. Originaliame darbe, kuriame buvo pasiūlytas reakcijos-difuzijos modelis, skaičiavimai buvo atliekami labai mažoje fizinėje erdvėje, kurioje telpa vos kelios mikrodalelės. Kadangi maišymo poveikiui didelę įtaką gali daryti erdvės, kurioje vyksta reakcija, dydis, šiame tyrime bus nagrinėjamos ir didesnės skaičiavimo sritys, siekiant geriau suprasti šį efektą.

\subsection{Darbo tikslas}

Šio darbo tikslas yra sudaryti kompiuterinę difuzijos-reakcijos sistemą, kuri modeliuotų kietafazę YAG reakciją bei ištirti maišymo modelio poveikį YAG kristalų sintezės spartai modeliuojant reakcija įvairaus dydžio erdvėse.

\subsection{Keliami uždaviniai ir laukiami rezultatai}
\subsubsection{Uždaviniai}
\begin{enumerate}
  
    \item Sudaryti kompiuterinius YAG reakcijos modelius remiantis baigtinių skirtumų skaitiniais metodais
    \item Patikrinti kompiuterinių modelių rezultatų korektiškumą bei palyginti modelių rezultatus tarpusavyje
    \item Apibrėžti medžiagų maišymo modelius
    \item Integruoti medžiagų maišymo modelius į skaitinius YAG reakcijos modelius
    \item Ištirti kompiuterinių modelių rezultatus įvairaus dydžio erdvėse
\end{enumerate}
\subsubsection{Laukiami rezultatai}
\begin{enumerate}
    \item Identifikuoti skaitiniai metodai, kurie bus naudojami YAG reakcijos modelio įgyvendinimui

    \item Įgyvendinti korektiškai veikiantys YAG reakcijos modeliai remiantis identifikuotais skaitiniais metodais

    \item Apibrėžti maišymo modeliai integruoti į kompiuterinius YAG reakcijos modelius

\end{enumerate}

\subsection{Tyrimo metodai}

\begin{enumerate}
    \item \textbf{Modeliavimo metodas}. Kompiuteriniai modeliai tyrimui bus konstruojami šiais baigtinių skirtumų metodais:
    \begin{itemize}
        \item Išreikštinis FTCS (\textit{angl. explicit forward time, centered space}) metodas
        \item Neišreikštinis kintamosios krypties metodas (\textit{angl. alternating direction implicit, ADI}) 
    \end{itemize}
    \item \textbf{Auksinio pjūvio paieškos algoritmas}. Ieškomas optimalus medžiagų išmaišymo laikas naudojant auksinio pjūvio paieškos algoritmą.
    \item \textbf{Lyginamoji analizė}. Lyginami skirtingais skaitiniais metodai įgyvendintų kompiuterinių YAG reakcijos modelių rezultatai. 
\end{enumerate}


\subsection{Numatomas darbo atlikimo procesas}

\begin{enumerate}
    \item Analizuojama literatūra 
    \item Sudaromi kompiuteriai YAG reakcijos modeliai naudojant python programavimo kalbą su \hbox{\textit{NumPy}} bei \textit{SciPy} paketais
    \item Tikrinamas kompiuterinių YAG reakcijos modelių rezultatų korektiškumas, lyginami skirtingų kompiuterinių modelių rezultatai. Duomenų vizualizavimui naudojamas \textit{Matplotlib} paketas
    \item Apibrėžiami maišymo modeliai
    \item Maišymo poveikio bei korektiškumo tyrimas kompiuterinių YAG reakcijos modelių rezultatuose 
    \item Tiriamas maišymo proceso poveikis įvairaus dydžio erdvėse
\end{enumerate}

% \subsection{Literatūros šaltinių apibūdinimas}

% \cite{ivanauskasComputationalModellingYAG2009,ivanauskasModellingSolidState2005,mackeviciusCloserLookComputer2012} - pasiūlyta matematinė difuzijos-reakcijos sistema, kuria modeliuojama YAG sintezės reakcija bei išmatuotos fizinės konstantos prie skirtingų aplinkos sąlygų, kurios leidžia kiekybiškai lyginti kompiuterinių modelių rezultatus. Šis modelis bus naudojamas įgyvendinti kompiuterinį YAG reakcijos modelį. \cite{pressNumericalRecipes3rd2007,levequeFiniteDifferenceMethods2007} - aprašo skaitinius metodus, kuriais gali būti sprendžiamos parabolinių dalinių diferencialinių lygčių sistemos. Šie šaltiniais bus naudojami literatūros analizei bei jais bus vadovaujamasi įgyvendintant kompiuterinius modelius.

\printbibliography[title = {Literatūra ir šaltiniai}]

\end{document}
